% Provisional empirical manuscript. Generic class; authors and contribution remain unconfirmed.
\documentclass[uplatex,a4paper,twocolumn,10pt]{jsarticle}
\usepackage[margin=20mm]{geometry}
\usepackage[dvipdfmx]{graphicx}
\usepackage{amsmath}
\usepackage{url}
% Generated from submission_metadata.json; no author names are inferred.
\newcommand{\LetterEnglishTitle}{Evaluating Information Retention and Execution Control in Stateful LLM Agents}
\newcommand{\LetterJapaneseKeywords}{ハーネスエンジニアリング，LLMエージェント，情報保持，実行制御，費用評価}
\newcommand{\LetterEnglishKeywords}{harness engineering, LLM agents, information retention, execution control, cost evaluation}

% Generated from closed P55 evidence; no model or task execution.
\newcommand{\InventoryAllCells}{%
\begin{tabular}{llrr}\hline 課題 & 条件・成否 & 仕事全体token & 継続生成\\\hline
写真整理 & F・未達 & 145,850 & 8\\
 & D・成功 & 120,935 & 6\\
 & B・未達 & 165,015 & 11\\
電気代精算 & F・未達 & 163,918 & 9\\
 & D・未達 & 161,225 & 9\\
 & B・未達 & 154,720 & 9\\
返金 & F・成功 & 85,697 & 2\\
 & D・成功 & 117,223 & 4\\
 & B・成功 & 117,305 & 4\\
曲一覧の書出 & F・未達 & 117,100 & 3\\
 & D・未達 & 156,265$+?$ & 7\\
 & B・未達 & 126,734 & 5\\
\hline\end{tabular}}
\newcommand{\InventoryTotals}{%
\begin{tabular}{lrrr}\hline 指標 & F & D & B\\\hline
成功／評価済み & 1/4 & 2/4 & 1/4\\
仕事全体・既知token & 512,565 & 555,648 & 563,774\\
継続・既知token & 302,185 & 345,268 & 353,394\\
使用量不明生成 & 0 & 1 & 0\\
前半＋継続生成 & 54 & 58 & 61\\
継続生成 & 22 & 26 & 29\\
前半＋継続API & 224 & 213 & 173\\
継続API & 173 & 162 & 122\\
復元API & 51 & 51 & 51\\
継続時間（秒） & 1,308.5 & 1,354.0 & 1,359.1\\
\hline\end{tabular}}
\newcommand{\InventoryFirstInputs}{%
\begin{tabular}{lrrr}\hline 課題 & F & D & B\\\hline
写真整理 & 8,789 & 7,342 & 6,848\\
電気代精算 & 11,267 & 11,009 & 10,202\\
返金 & 15,746 & 14,964 & 14,964\\
曲一覧の書出 & 13,771 & 12,142 & 9,473\\
\hline\end{tabular}}
\newcommand{\InventoryResultSummary}{全12予定条件のうち12件を評価し，0件が未評価，0件が未実行だった．継続費用は既知1,000,847 token，使用量不明は1生成だった．詳細省略Dは固定の採用基準を満たさず，品質を保った総費用削減として採用しない．Dは既知分だけでFを43,083 token上回るため，未知分を除いても削減にはならない．}
\newcommand{\InventoryContrasts}{全表示Fに対する全４課題の確定費用差は算出しない．一覧も省くBに対する全４課題の確定費用差は算出しない．}
\newcommand{\InventoryCaseEvidence}{写真整理のF/Dは，それぞれ42／42件のファイル移動要求を送ったが，成否は未達／成功だった．この対では全ての非supervisorアプリの保存ファイルが一致し，異なるアプリ状態はsupervisorだけだった．完了通知の要求はFで０，Dで１件であり，公式成否の差をファイル整理そのものの改善と呼ばない．隠れた個別テストの原因同定ではない．返金はF/D/Bの全条件が成功した．継続tokenはF/Dで32,209／63,735であり，初回の詳細省略量だけから終了までの費用を予測できない．}
\newcommand{\InventoryVariation}{過去P47と今回P55の全表示条件について，両方を評価できた比較のうち1課題で成否が異なった．これは事前に反復数を割り付けた追試ではなく，変更の原因や分散の推定には用いない．}

% P55/P71 all 24 planned display cells. Four distinct known tasks.
\newcommand{\PublicDisplayAllCells}{%
{\setlength{\tabcolsep}{2.5pt}
\begin{tabular}{llrrr}\hline 接続 & 課題 & F & D & B\\\hline
G & 写真 & 未達 145,850 & 成功 120,935 & 未達 165,015\\
 & 精算 & 未達 163,918 & 未達 161,225 & 未達 154,720\\
 & 返金 & 成功 85,697 & 成功 117,223 & 成功 117,305\\
 & 書出 & 未達 117,100 & 未達 156,265$+?$ & 未達 126,734\\
\hline
N & 写真 & 成功 156,754 & 未達 121,357$+?$ & 成功 129,261\\
 & 精算 & 未達 161,694 & 未達 162,725 & 未達 154,222\\
 & 返金 & 成功 85,975 & 成功 100,527 & 成功 100,527\\
 & 書出 & 未達 96,460 & 未達 138,798$+?$ & 未達 173,238\\
\hline
\end{tabular}}}
\newcommand{\PublicDisplayTotals}{%
\begin{tabular}{lrrr}\hline 指標 & F & D & B\\\hline
\multicolumn{4}{l}{文書取得必須（P55）}\\
成功／評価 & 1/4 & 2/4 & 1/4\\
仕事・既知token & 512,565 & 555,648$+?$ & 563,774\\
継続・既知token & 302,185 & 345,268$+?$ & 353,394\\
継続生成／不明 & 22/0 & 26/1 & 29/0\\
継続／復元API & 173/51 & 162/51 & 122/51\\
\hline
\multicolumn{4}{l}{公開名（P71）}\\
成功／評価 & 2/4 & 1/4 & 2/4\\
仕事・既知token & 500,883 & 523,407$+?$ & 557,248\\
継続・既知token & 290,503 & 313,027$+?$ & 346,868\\
継続生成／不明 & 21/0 & 25/2 & 27/0\\
継続／復元API & 204/51 & 138/51 & 148/51\\
\hline
\end{tabular}}
\newcommand{\PublicDisplayInitialSavings}{%
\begin{tabular}{lrrrr}\hline 接続 & 写真 & 精算 & 返金 & 曲書出し\\\hline
文書必須 & 1,447 & 258 & 782 & 1,629\\
公開名 & 1,447 & 258 & 782 & 1,629\\
\hline\end{tabular}}
\newcommand{\PublicDisplaySummary}{公開名接続では全12予定条件のうち12件を評価，0件が未評価，0件が未実行だった．継続費用は既知950,398 token，使用量不明は2生成だった．}
\newcommand{\PublicDisplayContrasts}{DとFの全４課題の確定費用差は算出しない．DとBの全４課題の確定費用差は算出しない．固定した開発上の採用条件は満たさなかった．}

\title{状態保持型エージェントのAPI説明省略における\newline 仕事全体の費用と終了手続きの診断}
\author{検討用原稿・著者情報未設定}
\date{2026年10月23日}
\begin{document}
\raggedbottom
\maketitle
\noindent\textbf{あらまし}\quad
\input{papers/202610_letter/cost_completion_abstract_ja.txt}

\noindent\textbf{キーワード}\quad\LetterJapaneseKeywords

\section{目的と検証する問い}
必要情報・処理と最終品質を保つ費用削減には，取得から検証までを含む評価が必要である．
短い入力でも再取得や生成が増えれば費用は減らず\cite{beyond,interaction}，
成功判定の差も実作業と終了手続きの違いを区別しない．
APIの一覧と引数の詳細を別々に省き，
\textbf{同じ途中状態からの表示省略が仕事全体費用と成否へどう影響し，
成否差を実作業の改善と解釈できるか}を調べる．
Python状態・スキーマ・原文参照を保持した既知４課題・２接続の全24表示条件と，
完了通知の介入，別条件の通常終了制御・推論保持を報告する．全結果と不明費用を残す．

\section{関連研究と本稿の範囲}
圧縮と成否・費用・停止理由の反復比較\cite{beyond}，再取得費用の定式化\cite{interaction}，
推論圧縮と後続の行動・全体費用の関係\cite{thinking}が既に調べられている．
短縮が総費用削減を保証しないという一般論は新規性としない．
CodeAct\cite{codeact}のコード行動，CaveAgent\cite{cave}の持続状態，
HD-Gist\cite{hdgist}のAPI引数・値の学習表現，Tool Search Tool\cite{toolsearch}の定義展開，
同一文書での段階的開示\cite{disclosure}も関連する．

同じ計算済み状態からの返却比較\cite{evidence}や，作用・方針・課題指標の分離\cite{matoi}も既出である．
ECT\cite{ect}は契約と記録を検証して終了し，AgentLoop\cite{agentloop}は要求充足・状態変化から反復を抑える．
検証費用のハーネス依存性も報告されている\cite{waste}．
本稿は取得済み説明の表示と完了通知を具体的に照合する．
既知要求の手書き検査を新しい終了方式や先行方式の再実装比較としない．

\section{費用の定義}
仕事全体の論理トークンを
\begin{equation}
 T_{\rm work}=T_{\rm prefix}+\sum_{j=1}^{n}(I_j+O_j)
\end{equation}
とする．共通前半の実測値を各条件へ一度加え，キャッシュ入力も$I_j$へ含める．
継続だけの値も示し，初回短縮量を全体の節約量へ，論理tokenを請求額・計算量・時間へ換算しない．
API，時間，実験用復元，評価，推論器のロード・常駐は別記し，重なる時間を足さない．
評価器単独の時間は表示比較では未計測で，稼働時間の残差からは求めない．

\section{表示比較の方法}
% Method description prepared during P71. No results are asserted here.
継続点の状態を会話履歴$H$，Python名前空間$P$，操作先環境$E$，
残予算$\mathcal B$，初期公開分を含む取得済み公開スキーマの集合$V$とする．
登録された全公開APIの集合を$\mathcal A$とし，$V\subseteq\mathcal A$である．
文書取得必須の接続は$V$，公開名接続は$\mathcal A$に属するAPIを呼び出せる．
後者でも，未取得のスキーマや具体的な引数値を追加表示するわけではない．
未登録名・未知の引数名はハーネスで拒否し，型・必須引数は元のAPI側で検査する．
認証，予算，完了後・作用不明後の停止も共通である．

全表示F，古い詳細のみ省略D，同じ詳細と古い一覧を省略Bを比較する．
保存した公開API返却と表示処理から出自を照合できる区間のみを変更し，
操作コード，リテラル，混合出力，エラー，最新応答を保持する．
D/Bで詳細の省略集合と印をそろえ，$P,E,V$と原文参照経路を維持する．
表示省略は保存情報の消去や呼出し権限の縮小とは異なる．

既知４課題の第８応答後から各条件を１回ずつ実行する．
文書取得必須G（P55）の結果を保存し，公開名N（P71）の比較を後に固定した．
P71では全３条件に同じ能力説明とツール説明を与え，前半の行動と返却は変更しない．
前半は元の文書取得必須条件で生成されたもので，そのときの拒否も履歴へ残す．
公開名接続を適用するのは継続部分であり，初期状態からの通常実行の比較ではない．
各比較内でD対FとD対Bを評価する．P55とP71の間では説明文と実行時点も異なるため，
その差を呼出し制約だけの因果効果とはしない．
両比較は同じ既知課題と前半を使い，独立した８課題の確認評価ではない．

各条件16生成，継続120,000 token，480秒を上限とし，全入力を計数して残予算内に出力枠を収める．
初回入力はキャッシュを再利用しない．MiniMax-M3のIQ4\_XS版とllama.cpp b11669を用い，
context 65,536，温度1.0，top-p 0.95，最大出力4,096とした．
seedは50から生成ごとに１増やし，順序は課題ごとにFDB，DBF，BFD，FBDとする．
API上限は前半込みで写真160，他80である．
使用量不明の生成は実行も再試行もせず当該条件を終了し，
作用・片付け・推論器の待機を確認できた場合だけ保存状態を評価して次へ進む．
不明費用，未評価，未実行を補完しない．

両比較でPython出力，公開返却，保存ファイルを継続前に照合するが，
ヒープ・乱数状態の完全複製や生成の決定性は主張しない．
返金D/Bは追加で省く一覧がなく基礎入力が同じである．
残時間通知と数値実装の揺らぎは残るため，そこでの差を一覧保持の効果とはしない．

% P75 scope clarification; no P76 result or quality effect is asserted.
会話履歴$H$は公開content・呼出し・観測であり，
全表示Fを含む全条件でモデルが返した推論は再投入しない．
したがって，表2・3の結果は推論保持を含む通常対照へ一般化できない．


\subsection{対象と入力の照合}
AppWorld\cite{appworld}の写真整理，電気代精算，返金，曲一覧の書出を用いる．
過去の失敗・入力計数を見た方式探索であり，未閲覧12シナリオは使わない．
D/Bの詳細省略は課題順に７・２・３・７返却で一致し，Bの追加一覧省略は１・１・０・３返却だった．
表1は共通の予算通知と省略印を含む初回短縮量である．

\begin{center}\begin{minipage}{\columnwidth}\centering
\small 表1\quad Dの初回入力短縮token（対F）\par\smallskip
\PublicDisplayInitialSavings
\end{minipage}\end{center}

\subsection{採用基準と欠測}
全12評価・全使用量既知・監査合格の下，DがF/B双方の成功を失わず，
合計$T_{\rm work}$がF未満・B以下なら局所採用候補とする．
一覧の価値には入力が異なるD成功・B失敗の対も要求し，同数成功を非劣性とはしない．
不明使用量は既知分と分け，ゼロや上限で補わず，確定費用差を求めない．
API作用・復元・実行・片付け・評価・推論器版の不一致では全体を停止する．
結果でseed・条件・予算を変えず，未実行の成否も補完しない．

\section{費用と成否の全結果}
両接続とも全12条件を評価し，未評価・未実行０だった（表2）．
継続費用はGで既知1,000,847 token＋不明１生成，Nで950,398 token＋不明２生成．
Dの既知費用はFよりGで43,083，Nで22,524 token多かった．
Nの写真はF/B成功・D未達で，Dの写真・書出しは生成中の時間上限により使用量不明となった．
未達で早く止まる条件を効率のよい成功とはしない．

\begin{center}\begin{minipage}{\columnwidth}\centering
\small 表2\quad 全24予定条件の成否・仕事token\par\smallskip
\PublicDisplayAllCells
\par\smallskip\footnotesize G：文書必須，N：公開名で継続．$+?$は別に使用量不明がある値．
\end{minipage}\end{center}

\subsection{入力短縮と仕事全体の費用}
不明生成のためD対F/Bの確定費用差は求めない．両比較ともDの既知費用がFを上回り，採用条件は未達だった．
表3に合計を示す．GのDは初回入力がFより短いが，継続生成は22回から26回へ増えた．
返金は両接続のF/D/Bで成功したが，Dの全体費用はFより大きかった．

\begin{center}\begin{minipage}{\columnwidth}\centering
\small 表3\quad 全４課題合計の費用と実行量\par\smallskip
\PublicDisplayTotals
\end{minipage}\end{center}

継続APIは実際の送信数で，拒否や原文参照は別記する．回数減だけを不要取得削減とはしない．
推論器全稼働はGで4,112.859秒，Nで3,724.097秒．Nの継続490・復元153 APIのうち，
説明公開前の呼出しは写真Bの公開時計１件だった．

\section{成否差と終了手続きの診断}
Gの写真F/Dは各42移動で非supervisor保存ファイルが一致したが，公式成否は未達／成功，通知は０／１件だった．まずGの写真F/Bを別環境で再現し，元Dと同じ空引数の完了通知なし／１回追加の４条件を固定した．
次にNの全12条件のうち，成否によらず完了API未送信の全７状態（既知３課題）を選び，同じ２条件の計14ケースを固定した．Nの通知は元写真Bで観測した\texttt{status="success"}だけとした．全ケースで送信済み公開APIの返却と介入前の全13保存ファイルを元状態と照合し，無介入は元の成否を再現した．
\begin{center}\small 表4\quad 完了通知だけの介入（全18ケース）\par\smallskip
\begin{tabular}{lll}\hline 再現した元状態 & 通知なし & 通知１回\\\hline
G 写真 F & 未達 & 成功\\
G 写真 B & 未達 & 未達\\
N 写真 D & 未達 & 成功\\
N 精算 F & 未達 & 未達\\
N 精算 D & 未達 & 未達\\
N 精算 B & 未達 & 未達\\
N 書出し F & 未達 & 未達\\
N 書出し D & 未達 & 未達\\
N 書出し B & 未達 & 未達\\
\hline\end{tabular}\end{center}
Gの写真FとNの写真Dだけが未達から成功へ変わり，それぞれ元GのD，元NのBと最終13保存ファイルも一致した．非supervisor保存ファイルは全ケースで介入前後不変だった．残り７状態は通知後も未達で，完了宣言を一律に足す根拠にはならない．
Gの診断は346 API・４評価・23.927秒（評価8.541秒），Nは702再送＋７通知＝709 API・14評価・64.318秒（評価19.503秒）だった．両診断とも生成０で，元の成否と不明費用は変更しない．Nの不明２生成は最後の未実行応答に限り，全API作用が既知であることを別に照合した．
同じ既知写真の事後診断であり，独立した別課題の追認ではない．Gの写真では分類に必要な現在年の成功した公開時計応答がなく，通知による公式判定の変化は公開証拠だけで要求を全て検証できたことを意味しない．


\section{公開証拠による通常終了制御の追加点検}
同じ既知写真を初期状態から文書必須・公開名の各接続で実行した別の対照を用いた．
モデル・seed・案内・予算をそろえ，応答内の全プログラム実行後に公開記録を検査した．
上限は各80生成・450,000 token・1,800秒・160 API．
元集合・移動前日時・現在年・移動先・名前保持が支持されれば完了APIを１回送る．
不明時は通常継続し，追加案内・未来応答・評価器は使わない．
% Generated from the closed P66 audit by build_cost_completion_paper.py.
\begin{center}\begin{minipage}{\columnwidth}\centering
\small 表5\quad 別条件の写真課題：全体token\par\smallskip
\begin{tabular}{lrrr}\hline
接続 & 通常終了 & 終了制御 & 減少率\\\hline
文書取得必須 & 130,439 & 91,116 & 30.1\%\\
公開名 & 151,664 & 112,233 & 26.0\%\\
\hline\end{tabular}
\end{minipage}\end{center}

全４実行が成功し，制御では各２生成・４ APIが少なかった．
各対の最終13保存ファイルと制御時点までの公開API記録は一致したが，
ID差による要求・モデルメッセージ全文一致の事前基準は不成立だった．

制御２条件の全25生成は203,349 token，199 API，２評価，使用量不明０だった．
25回の検査と記録ダイジェスト計算は合計0.028224秒で，全体時間にも含めた．
減少78,754 tokenのうち66,720はcached inputであり，金銭・時間の同率削減ではない．
既知書出しでは通常ツール対照と終了制御を各１回，API上限のみ80として比較し，
正しいCSV作成後の退会までを公開記録に要求した．
両方８生成・219,126 token・80 APIで未達，書出し・退会０回．
８検査0.004080秒で通知・削減０，使用量不明０だった．最終13保存ファイルは一致し，ID差の全文不一致は残った．
このID差を別途点検し，固定runtimeで保存66入力と３対照を再構成した．元batchは原文と正規化後のテンプレートの末尾改行１byte差で停止したため，保存出力を事後に読み取った．写真13・12対，書出し８対の全33対で文章・全token列・生成制約が一致し，元66入力の実測数とも一致した．未変更とID変更の対照は一致，本文変更では文章・token列が変化した．69入力の再構成・計数を含む処理5.543秒，事前照合1.859秒，事後読取0.254秒で，新規生成・API・評価は０だった．過去のGPU・乱数器状態は復元せず，元の全文不一致判定は変更しない．

これらはAPIの意味を信頼した既知文型の通常制御であり，表2・3へ合算しない．
実装・固定条件・全件集計を公開する\footnote{\url{https://github.com/sophiria1225/AI-}}．

% All eight P76 conditions; generated only from the closed audited batch.
\section{推論履歴を保持する通常対照}
同じ既知４課題の８応答後から，公開履歴のみ（P）と返却推論も全て保持（R）を各１回比較した．
両条件は公開履歴・持続Python・公開名接続・予算を共通にし，過去と以後の推論の投入だけを変えた．
各16生成・120,000 token・480秒で，最初の生成はcacheを使わない．
共通前半は推論を再投入せず生成しており，最初からの方針比較ではない．
\begin{center}\begin{minipage}{\columnwidth}\centering
\small 推論保持比較の全８条件：成否・仕事token\par\smallskip
\begin{tabular}{lrr}\hline 課題 & P：公開履歴のみ & R：推論も保持\\\hline
写真 & 成功 156,754 & 成功 118,184\\
精算 & 未達 160,878 & 未達 164,322\\
返金 & 成功 85,720 & 成功 88,001\\
書出 & 未達 134,971$+?$ & 未達 100,425$+?$\\
\hline
成功／評価 & 2/4 & 2/4\\
既知token合計 & 538,323$+?$ & 470,932$+?$\\
\hline\end{tabular}
\end{minipage}\end{center}
Rの初回入力は課題順に291・189・1,116・361 token増えた．
全８予定条件のうち8条件を評価し，使用量不明は2生成だった．欠測を補完せず，採用判定に用いる確定費用差は求めない．
課題別のR−Pの仕事token差は，写真 -38,570／精算 +3,444／返金 +2,281／書出し 未確定だった．
全８評価・全使用量既知・少なくとも１対成功・成功喪失０・全体費用減少という局所採用条件は満たさなかった．
新規継続は43生成・既知588,495 token＋使用量不明2生成，継続306・復元102 APIだった．
推論器全稼働2,662.086秒に復元・評価も含め，重複加算しない．
保持は通常の比較条件であり，新規手法や一般的な非劣性の実証とはしない．表2・3の過去条件は変更しない．


\section{解釈の限界}
同じ既知４課題・単一モデル・各条件１回の開発比較であり，
診断18ケースも既知３課題９状態への事後介入である．
seedと初回cacheをそろえても後続経路の再現性は保証しない．
返金D/Bは基礎入力が同じでも後続の残時間通知等が異なり，差を一覧の効果に帰属させない．

Bでも過去コード等にAPI名が残る（写真8/26名，精算10/26名）．
名前の理解や情報の完全消去を測るものではなく，省略対象は出自の確かな古い説明に限る．

\section{結論}
両接続で，Dの初回入力は短くても既知の仕事全体費用はFを上回った．
写真の成否差と終了通知の関係，通常終了制御の局所改善と未達を示した．
推論保持でも写真の費用は減ったが，失敗・不明使用量が残り，全体の採用条件は未達だった．
既知４課題の局所結果であり，必要情報と品質を守る制御には，
仕事全体費用と作業・終了手続きへの作用を対応させた検証が必要である．

% Full-author, sentence-case bibliography; metadata and byte hash are checked separately.
\begin{thebibliography}{14}\small
\bibitem{beyond} R. Satish, P. Sinha, A. Kawada, and N.J. Yadwadkar, ``Beyond token savings: a systematic study of context compression in LLM agents,'' arXiv:2609.32961v1, Sept. 2026.
\bibitem{interaction} S. Liu, ``What does context compression cost an agent? Interaction costs unrevealed by task-completion metrics,'' arXiv:2608.16370v1, Aug. 2026.
\bibitem{thinking} M. Wang, F. Luo, B. Wang, G. Yao, Y. Guo, C. Ning, H. Chen, Y. Ma, and J. Han, ``When can agents forget their reasoning? ICLR for long-horizon agent context compression,'' arXiv:2609.29875v1, Sept. 2026.
\bibitem{codeact} X. Wang, Y. Chen, L. Yuan, Y. Zhang, Y. Li, H. Peng, and H. Ji, ``Executable code actions elicit better LLM agents,'' Proc. ICML, PMLR, vol.235, pp.50208--50232, Vienna, Austria, July 2024.
\bibitem{cave} M. Ran, Z. Wan, C. Lin, Y. Zhang, H. Xin, H. Fan, Y. Xu, B. Luo, Y. Zhou, W. Zhao, L. Yang, L. Feng, F. Yang, J. Wu, Y. Huang, C. Ma, Y. Huang, D. Jiang, J. Deng, S. Han, Y. You, B. An, Y. Guo, and J. Song, ``CaveAgent: transforming LLMs into stateful runtime operators,'' arXiv:2601.01569v4, June 2026.
\bibitem{hdgist} Y. Jiang, M. Del Vecchio, M. Bansal, and A. Johannsen, ``Hierarchical and dynamic prompt compression for efficient zero-shot API usage,'' Findings of EACL, pp.2162--2174, St. Julian's, Malta, March 2024.
\bibitem{toolsearch} Anthropic, ``Introducing advanced tool use on the Claude Developer Platform,'' Nov. 24, 2025. \url{https://www.anthropic.com/research/advanced-tool-use}, accessed Oct. 23, 2026.
\bibitem{disclosure} Y. He, Y. Zhao, J. Wang, and H. Chen, ``Is progressive disclosure all you need for long-context agents?,'' arXiv:2607.17598v1, July 2026.
\bibitem{evidence} K. Zhang, S. Chundur, J.P. Owusu, M.J. Abdolhosseini Qomi, and M. Raissi, ``Same simulation, different evidence: where retrieved tool output costs LLM agents accuracy,'' arXiv:2610.26836v1, Oct. 2026.
\bibitem{matoi} Z. Wang, ``Matoi: selective revalidation of consequence evidence for agent authorization,'' arXiv:2610.21394v1, Oct. 2026.
\bibitem{ect} J. Liu, ``When may an agent stop? Evidence-carrying termination for tool-using LLMs,'' arXiv:2608.23623v1, Aug. 2026.
\bibitem{agentloop} W. Zheng, M. Xu, K. Hu, K. Ye, and C. Xu, ``AgentLoop: runtime control of slot-closed execution loops for tool-augmented LLM agents,'' arXiv:2609.33315v1, Sept. 2026.
\bibitem{waste} S. Weinberger and A. Hozez, ``Prompt-induced waste in coding agents: reasoning, effort, harness design, and end-to-end cost,'' arXiv:2608.01347v6, Sept. 2026.
\bibitem{appworld} H. Trivedi, T. Khot, M. Hartmann, R. Manku, V. Dong, E. Li, S. Gupta, A. Sabharwal, and N. Balasubramanian, ``AppWorld: a controllable world of apps and people for benchmarking interactive coding agents,'' Proc. ACL, vol.1, pp.16022--16076, Bangkok, Thailand, Aug. 2024.
\end{thebibliography}

% English abstract and keywords, outside the standard body page count.
\label{LastBodyPage}\clearpage\onecolumn
\section*{English abstract and keywords}
{\large\bfseries \LetterEnglishTitle\par}\bigskip
\noindent\textbf{Abstract}\quad 
\input{papers/202610_letter/cost_completion_abstract_en.txt}

\noindent\textbf{Keywords}\quad\LetterEnglishKeywords

\end{document}
